\scanpage{54}
\mspage{13}
\deleted{\unsure{Prouvons|Prenons|Proposition}.}

\uline{Corollaire}\quad $\operatorname{Card} W=72\mathbin.6!=2^7\mathbin.3^4\mathbin.5$

En effet, $W$ est \uline{simplement} \uline{transitif} sur $V$ ens des bases ordonnées, qui \unsure{sont au nb|est au nb|a un card} de $72\mathbin.6!$. (car il y a $72$ parties basiques i.e. $72$ racines \ldots)

\uline{Th}\quad $(E(\Delta),R(\Delta))$ est un \deleted{\unsure{ens de|un}} \inserted{système de} racines de type $E_6$, \inserted{et $E(\Delta)$ engendré par $R$.} et si $(\zeta_1,\ldots\zeta_6)$ est une base ordonnée de $\Delta$, le système de racines
\[
\begin{aligned}
&\deleted{\alpha_1=\zeta_1\unsure{-}\zeta_2}\\
&\deleted{\unsure{\alpha}}r_1=r_{12}=\zeta_1-\zeta_2\\
&\deleted{\unsure{\alpha}}r_2=r_{23}=\zeta_2-\zeta_3\\
&\ldots\\
&r_{\deleted{6}5}=r_{56}=\zeta_5-\zeta_6\\
&r_6=r_{123}=\deleted{\unsure{2\eta-\zeta=2\eta-\textstyle\sum\zeta_i}}\ \eta-\zeta_1-\zeta_2-\zeta_3
\end{aligned}
\]
est un système de racines simples, suivant le diagramme
\diagram{\begin{tikzcd}[ampersand replacement=\&,column sep=1.1em,row sep=.65em]
\overset{r_1}{\bullet}\arrow[r,dash]\&\overset{r_2}{\bullet}\arrow[r,dash]\&\overset{r_3}{\bullet}\arrow[r,dash]\arrow[d,dash]\&\overset{r_4}{\bullet}\arrow[r,dash]\&\overset{r_5}{\bullet}\\
\&\&\underset{r_6}{\bullet}
\end{tikzcd}}

\inserted{Pour $r\in R$,} La symétrie $s_r$ restreinte à \inserted{$E={}$}$E(\Delta)$ est la sym. associée au syst. de racines et à $r$, et $W$ (identifié à un groupe d'autom de $E$) est \inserted{donc} le groupe de Weyl du syst. de racines. La forme bil. \inserted{sym.} $\varphi=\overline\varphi\mid E$ est caractérisée par la propriété d'être invariante \inserted{par $W$} et de satisfaire $r.r=-2$ $\forall r\in R$.
% archive sheet number: 53
